\documentclass[11pt]{article}
\usepackage{amsmath,amssymb,amsthm}
\newtheorem{theorem}{Theorem}
\newtheorem{lemma}[theorem]{Lemma}
\newtheorem{proposition}[theorem]{Proposition}
\title{Permutation Parity Fingerprints and the Missing Minor Reader}
\author{}
\date{November 24, 2026}
\begin{document}
\maketitle

\paragraph{Status and scope.}
The combinatorial theorem below is a complete written proof. It supplies
an explicit query family with order $H\log H$ independent binary
readouts from permutations of $H$ symbols. It does not implement those
queries by attention heads. The exact transfer gap is recorded at the
end. No model definition, Lean declaration, or asymptotic head bound is
changed by this note.

Fix $H\geq2$ and strictly increasing real numbers
$\lambda_1<\cdots<\lambda_H$. For a subset $S\subseteq\{1,\ldots,H\}$
and a permutation $\pi$, define its induced-order parity by
\[
 \chi_S(\pi)=\operatorname{sgn}
 \prod_{\substack{h<k\\h,k\in S}}
       (\lambda_{\pi(k)}-\lambda_{\pi(h)}).
 \tag{1}
\]
An empty product is one. Every factor is nonzero, so all these signs
are strict. The query value depends only on the order of the tokens,
not their spacing.

\begin{lemma}[Separation probability]
For distinct permutations $\pi,\sigma$ and a uniformly random subset
$S\subseteq\{1,\ldots,H\}$,
\[
 \mathbb P\{\chi_S(\pi)\ne\chi_S(\sigma)\}\geq\frac14.
 \tag{2}
\]
\end{lemma}
\begin{proof}
Write $s_h=\mathbf1_{h\in S}$. Over the field $\mathbb F_2$, the
parity of the ratio of the signs in (1) is the multilinear polynomial
\[
 Q(s)=\sum_{\substack{h<k:\\
  \mathbf1_{\pi(h)>\pi(k)}\ne
  \mathbf1_{\sigma(h)>\sigma(k)}}} s_hs_k.
 \tag{3}
\]
The sum is nonempty: the pairwise order comparisons uniquely determine
the rank of each entry and hence the permutation. Thus $Q$ is a
nonzero multilinear polynomial of degree two.

For completeness, every nonzero multilinear polynomial of degree at
most $d$ over $\mathbb F_2$ takes the value one on at least a fraction
$2^{-d}$ of the Boolean cube. The degree-zero assertion is immediate.
For positive degree, choose a variable occurring in a monomial of
maximum degree and write the polynomial as $x_iA+B$, with $A\ne0$
of degree at most $d-1$. Conditional on the other coordinates and on
$A=1$, exactly one of the two choices of $x_i$ makes the polynomial
one. Induction therefore gives probability at least
$\tfrac12 2^{-(d-1)}$. Applying this with $d=2$ proves (2).
\end{proof}

\begin{theorem}[A small separating query family]
There exist
\[
 M=\left\lceil8\log_2(H!)\right\rceil
 \tag{4}
\]
subsets $S_1,\ldots,S_M$ such that the binary vectors
$(\chi_{S_j}(\pi))_{j=1}^M$ are distinct for all $H!$ permutations.
Repeated subsets, if any, can be discarded.
\end{theorem}
\begin{proof}
Choose the subsets independently and uniformly. A fixed pair of
permutations has identical query vectors with probability at most
$(3/4)^M$. A union bound makes the probability of any collision at
most
\[
 {H!\choose2}(3/4)^M<\frac{(H!)^2}{2}
                  (3/4)^{8\log_2(H!)}<\frac12.
\]
For the last inequality, $(3/4)^4=81/256<1/2$, so
$(3/4)^8<1/4$. A collision-free realization consequently exists.
\end{proof}

\begin{theorem}[Many independent strict parity queries]
There is a family of more than $\log_2(H!)/10$ subsets whose parity
queries jointly realize every binary word as $\pi$ varies. In
particular, for $H\geq4$ its size is greater than
$H\log_2 H/40$.
\end{theorem}
\begin{proof}
Put $N=H!$ and $L=\log_2N\geq1$. Take the separating family above,
which has $M\leq9L$ coordinates, and let $d$ be the largest number
of its coordinates shattered by the $N$ distinct sign vectors.
The Sauer counting inequality gives
\[
 N\leq\sum_{j=0}^d{M\choose j}.
 \tag{5}
\]
Here is the relevant counting argument. For a family of binary
vectors, project away the last coordinate. Its size equals the
number of projected vectors plus the number of projected vectors
having both last-coordinate extensions. The first projected family
has shattering dimension at most $d$; the second has dimension at
most $d-1$, since every shattered set there remains shattered after
adjoining the last coordinate. Induction on $M$, using Pascal's
identity, proves (5).

Since $N>1$, we have $d\geq1$. If $d\leq L/10$, then $d\leq M$
and the usual binomial estimate gives
\[
 N\leq(eM/d)^d\leq(90e)^{L/10}<2^L=N,
 \tag{6}
\]
a contradiction. For clarity, the first estimate follows by putting
$t=d/M\leq1$ and bounding the sum in (5) by
$t^{-d}(1+t)^M\leq(eM/d)^d$. The second follows because
$x\log(eM/x)$ is increasing on $0<x\leq M$, and because $M\leq9L$.
The strict last inequality follows already from $e<3$ and
$270<2^{10}$. Thus $d>L/10$.

Finally $H!\geq(H/2)^{H/2}$, by retaining the largest
$\lceil H/2\rceil$ factors of the factorial. For $H\geq4$ this
implies
$L\geq(H/2)(\log_2H-1)\geq(H/4)\log_2H$.
\end{proof}

\section{What must still be transferred to attention}
Equation (1) is the sign of a Vandermonde minor. If
$S=\{h_1<\cdots<h_s\}$, it equals the sign of
\[
 \det\bigl[\lambda_{\pi(h_j)}^{i-1}\bigr]_{i,j=1}^s.
 \tag{7}
\]
For strictly ordered exterior Cauchy row arguments, the corresponding
Cauchy minor has this same sign up to a known fixed sign. Thus an
original-head construction that simultaneously reads these variable
subset-minor signs with $O(H)$ heads, under the compact whole-cube
denominator budget, would yield the desired order $H\log H$ local
capacity. The preceding theorem does not supply that construction.

In particular, solving a fixed system of $H$ zero-anchor equations
produces entries of one inverse matrix. A scalar added query gives
a linear combination of cofactors divided by its determinant.
The matrix determinant lemma expresses this as one rank-one update.
It does not identify this quantity with the arbitrary principal
minors required by (7). Jacobi's complementary-minor identity uses
determinants of submatrices of the inverse, rather than a single
entry or a single linear query, so it does not repair this gap.

There is also an elementary obstruction to a proposed common-probe
shortcut. Suppose the solved residues are token-only, $c_h=c(\lambda_{\pi(h)})$,
and a query changes a common probe on all slots in $S$, leaving the
other slots at another common probe. Its score is a sum of one
token function over $S$ and another over its complement. Swapping
two tokens within $S$ leaves that score unchanged, whereas it
reverses (1). This fails already for $|S|=2$. Slot-dependent or
genuinely interacting residues are necessary for this parity reader.

\paragraph{Current disposition.}
The common-probe and one-cofactor implementations above are blocked.
No claim of impossibility is made for general admissible ratio sums.
A further construction must explicitly implement the subset-minor
signs, or provide another fingerprint family with the separation
property (2), without assuming a determinant evaluator as a head.
\end{document}
